\chapter{合成有机高分子化合物}
我们已经学过的淀粉、纤维素、蛋白质等都是天然的有机高分子化合物。
学过的聚乙烯、聚氯乙烯、酚醛树脂等都是人工合成的有机高分子化合物。
在这章里，主要叙述合成有机高分子化合物的结构、性质、合成和用途等，但也涉及到某些天然有机高分子化合物，如天然橡胶。

\section{概述}
\subsection{有机高分子的式量}
我们知道，烃、醇、醛、羧酸、酯、单糖、低聚糖的式量都比较低，如蔗糖的式量是 342，硬脂酸甘油酯的式量是 890。
它们的式量很少上千，可以称为低式量的化合物或简称小分子。
相反，淀粉的式量从几万到几十万，纤维素的式量也达几十万，蛋白质的式量从几万到几百万或更高。
核蛋白的式量迄今所知，是最高的，有的达几千万。
由乙烯合成的聚乙烯，由氯乙烯合成的聚氯乙烯，它们的式量一般达几万到几十万。
淀粉、纤维素、蛋白质、聚氯乙烯、聚乙烯、酚醛树脂等物质都属于高分子化合物。
高分子化合物简称高分子。

高分子虽然式量很大，但结构往往比较简单，它们是由结构单元重复构成的。
什么是高分子化合物的结构单元呢？
我们知道，纤维素的分子（\ce{(C6H10O5)_$n$}）是由成千个 \ce{C6H10O5} 连接起来的。
聚乙烯分子（\ce{(C2H4)_$n$}）是由成千成万个 \ce{C2H4} 连接起来的。
纤维素的结构单元是 \ce{C6H10O5}，聚乙烯的结构单元是 \ce{C2H4}。
高分子里的重复结构单元就是\Concept{链节}。
$n$ 表示每个高分子里链节的重复次数，叫\Concept{聚合度}。
聚乙烯是由乙烯聚合而成的，乙烯就是聚乙烯的\Concept{单体}。
高分子里的链节是由单体在反应中形成的。

高分子从单个分子来说，它有一定的聚合度，也就是说，$n$ 是某一个整数，所以它的式量是确定的。
但对于一块高分子材料来说，它是许多聚合度相同或不相同的高分子聚集起来的。
因此，我们从实验测得的某种高分子的式量只能是一种平均式量。
高分子的这种性质是跟小分子不同的。

\subsection{有机高分子的结构}
高分子区别于小分子还有结构上的特点。
单个高分子是由一个个链节连接起来的，成千上万的链节常常连成一条长链。
高分子最普通最重要的结构是长链状的。
这就是人们发现的高分子的线型结构（\cref{fig:5-1a}）。
高分子链中，原子跟原子或链节跟链节都是以共价键相结合的。
如聚乙烯、聚氯乙烯的长链都是由 \ce{C-C} 键连接的。
淀粉和纤维素的长链是由 \ce{C-C} 键和 \ce{C-O} 键相连接的。
蛋白质的长链是由 \ce{C-C} 键和 \ce{C-N} 键相连接的。
这些键都是单键，单键能够自由转动，所以高分子链里的 \ce{C-C} 键能够旋转。
高分子链上单键的旋转往往引起链节，更多地是多个链节（即链段）的旋转。
我们可以想象，一条能够旋转的长链，没有外力拉它，决不会是直线形的，高分子链确乎是蜷曲的长链。
\begin{figure}
  \begin{minipage}[b]{0.32\linewidth}\centering
    \includegraphics{5-1a.pdf}
    \subcaption{不带支链的线型结构}\label{fig:5-1a}
  \end{minipage}
  \begin{minipage}[b]{0.32\linewidth}\centering
    \includegraphics{5-1b.pdf}
    \subcaption{带支链的线型结构}\label{fig:5-1b}
  \end{minipage}
  \begin{minipage}[b]{0.32\linewidth}\centering
    \includegraphics{5-1c.pdf}
    \subcaption{交联的体型（网状）结构}\label{fig:5-1c}
  \end{minipage}
  \caption{高分子结构型式示意图}\label{fig:5-1}
\end{figure}

当许多高分子聚集在一起的时候，情况又是怎样呢？

我们知道，分子跟分子有着相互作用的力，这就是分子间力，分子间力的能量比化学键的键能小得多。
但对于高分子来说，情况就不一样了。
高分子是长链，链节是成千上万的，分子跟分子是缠在一起的，分子间接触的地方非常之多，也就是说许许多多处有分子间力。
所以高分子材料的强度是由高分子链中的化学键和分子间力表现出来的，有时更重要的是分子间力。
如果高分子的式量大，链也长，那么高分子间的分子间力也相应地增强了。
所以物质的式量必须达到一定的程度才能显出高分子材料的性质来。

线型结构的高分子有些是带支链的（\cref{fig:5-1b}），如支链淀粉是带支链的。
带支链的高分子仍然是一个个的单独的高分子。

高分子链上如果还有能起反应的官能团，当它跟别的单体或别的物质起反应的时候，高分子链之间将形成化学键，产生一些交联（\cref{fig:5-1c}），形成体型（网状）结构，下面要学到的硫化橡胶就是这样的例子。

如上面所说的，根据结构，高分子通常分为线型的和体型的。
线型的，包括有支链的；体型的，即网状的，有交联多的，也有交联少的。

\subsection{有机高分子的性质}
有机高分子由于它们的式量大和一定的结构而具有跟小分子物质不同的特殊的性质。
人们根据这些性质来应用它们。

根据使用的观点，通常把有机高分子材料分成塑料、纤维和橡胶等几大类。
它们都各有特殊的性质。
随着高分子科学和技术的发展，人们有可能有目的地把它们合成出来。
高分子材料的性质是多方面的，这里只介绍跟结构关系明显的、跟生产和使用关系密切的一些性质，即高分子材料在溶剂中的溶解、高分子材料在不同温度下的性能以及其它重要性质。

\subsubsection{高分子材料在溶剂中的溶解}
高分子材料在溶剂（通常是有机溶剂）里的溶解过程跟小分子的有所不同。
对线型高分子来说，溶解的第一步是溶剂分子渗透入缠在一起的线型高分子之间，使高分子材料胀大。
第二步是溶剂分子把高分子包围起来，使高分子一个个分离开来，能够自由移动。
线型高分子溶解在有机溶剂里成为高分子溶液。
溶液里高分子的大小达到胶体微粒的大小，成为溶胶，但高分子溶液里分散着的又是单个的高分子。

\begin{Experiment}
把聚苯乙烯树脂刮成粉末，取粉末 \qty{0.5}{g}，放在一个试管里，加入 \qty{10}{ml} 苯，观察溶解的情况。
\end{Experiment}

\begin{Experiment}\label{expr:5-2}
  把 \qty{2}{g} 苯酚放在一个试管里，加入 \qty{1.5}{ml} 三氯甲烷，把试管放在水浴中加热到苯酚完全溶解。加入 \qty{1}{g} 聚酰胺 6（锦纶）的丝，观察溶解的情况。
\end{Experiment}

\begin{Experiment}\label{expr:5-3}
取聚甲基丙烯酸甲酯（即有机玻璃）粉末 \qty{0.5}{g} 放在试管里，加入 \qty{9}{ml} 三氯甲烷。观察溶解的情况。
\end{Experiment}

从上面的实验，可以观察到，这些线型高分子是能溶解在适当的溶剂里的。
但是高分子的溶解过程比小分子缓慢。

高分子溶液对于研究高分子的式量、结构、性质和加工等是很重要的。

\begin{Experiment}
从废轮胎上刮下一些橡胶粉末，取 \qty{0.5}{g} 粉末放在试管里，加入 \qty{5}{ml} 汽油（或苯）。观察发生的现象，它能不能溶解。
\end{Experiment}
从上面实验可以看出，交联的体型高分子难于溶解，但能有一定程度的胀大。

\subsubsection{高分子材料在不同温度下的性能}
\begin{Experiment}
在三个试管里分别放聚氯乙烯、聚乙烯、聚丙烯的塑料碎片各 \qty{3}{g} 左右，用酒精灯缓缓加热。观察软化和熔化的情况。等熔化后即停止加热以防分解。等待冷却后又变固体。再加热又熔化。  
\end{Experiment}

可以观察到上述各线型高分子受热到一定温度范围，开始软化，直到熔化成流动的液体。
但不象固态小分子物质那样有确定的熔点。
冷却变固体后，加热又熔化。
所以线型高分子具有热塑性。
就是加热软化后，可以加工成各种形状。

当固态的线型高分子受热熔化后，分子里的链节、链段能够移动、旋转。
受外力作用时，整个分子链互相滑动，发生形状的变化，除去外力也不恢复原状。

随着温度下降，分子的动能减少，整个的分子链已不会相互滑动，但链段还能作移动旋转。
当受外力时，能够使高分子链的一部分由蜷曲到伸展；失去外力又恢复到蜷曲的原状。

当温度再下降的时候，分子的动能继续减少，无论分子活动或链段的旋转都已冻结，只有分子在它的位置上作振动了。

体型高分子受热不会熔化，这因为网状结构中的高分子链已有化学键互相交联，限制了高分子链移动的缘故。
当温度更高时，体型高分子的化学键出现断裂，高分子结构就被破坏。
所以体型高分子具有热固性，即一经加工成型后，就不再受热熔化。

但是，体型高分子材料和线型高分子材料的性质差别也是相对的，支链很多的线型高分子的性质接近于交联程度低的体型高分子的性质。

\subsubsection{高分子材料的强度、塑性、电绝缘性等}
有许多高分子材料的强度比较大。
如把 \qty{10}{kg} 高分子材料或金属材料做成 \qty{100}{m} 长的绳子，在高楼上吊重物。
那么，锦纶能吊 \qty{15500}{kg}，涤纶能吊 \qty{12000}{kg}，金属钛能吊 \qty{7700}{kg}，碳钢能吊 \qty{6500}{kg}。
高分子材料的密度比金属材料小，密度常在 \qtyrange{0.9}{1.5}{g/cm^3} 之间。
如果用同样截面积的绳子来比较所能承受的拉力，那么碳钢是 \qty{14500}{kgf/cm^2}，锦纶是 \qty{10500}{kgf/cm^2}。
所以对相同质量的材料来说，高分子材料一般具有强度大的性能。

在一定条件下，高分子材料具有好的可塑性，便于加工，能被拉成丝、吹成薄膜、能用模具压成各种所需要的形状等。
如用高分子材料吹成的薄膜，用于农业、工业、日常生活等；拉成的各种各样的丝，织成了布，制成了渔网，等等。

高分子链里的原子是以共价键结合的，一般不具有自由电子，不易导电，高分子材料常是很好的电绝缘材料，用于包裹电缆、电线，制成各种电器设备的零件，等等。

有些高分子是饱和烃的长链，具有饱和烃的稳定性，能耐化学腐蚀。
有的高分子链上有苯环结构，比较耐热，熔化温度高。
高分子材料还比较耐磨、不透水、或比较耐油等。

高分子材料虽然有许多优良特性，但也有一些缺点，一般说来，就是不耐高温，易燃烧，易老化等。
所谓老化就是高分子材料在加工或使用过程中，受到光、热、空气、潮湿、腐蚀性气体等综合因素的影响，使材料逐步失去原有的优良性能，以致最后不能使用。
通过改善高分子化合物的结构、改进它们的聚合和加工工艺以及注意使用的环境和条件，可以提高高分子的性能，减少或延迟老化等，这些都是对高分子材料的重要研究课题。

\begin{Practice}[习题]
\begin{question}
  \item 线型高分子和体型高分子在结构上、性质上，有哪些区别？
  \item 某种聚氯乙烯的聚合度是 2000，计算它的平均式量。（提示: 高分子的式量 $=$ 聚合度 $\times$ 链节的量。）
  \item 塑料口袋的封口是根据什么原理操作的？
  \item \cref{expr:5-2} 配制的高分子溶液可以用来修补稍有损坏的锦纶织品。\cref{expr:5-3} 配制的高分子溶液可以用来修补损坏的有机玻璃制品。可以在教师指导下进行试验。
\end{question}
\end{Practice}

\section{加聚反应和缩聚反应}
合成高分子的少数品种早在十九世纪就被人们制备出来了。
在二十世纪初还只有少量的生产。
由于性能好，很实用，逐渐受到重视。
但缺乏理论的指导，没有找到有效的生产方法。
到了三十年代，高分子结构理论初步确立，反应原理逐步阐明，合成高分子的许多品种迅速发展起来。
合成高分子的原料极其丰富多样，资源很广。
它们可以来自煤、天然气、石油、农副产品。
石油化工的迅速发展为高分子合成提供品种更多、数量更大的合成原料，如乙烯、丙烯、丁二烯、苯、甲苯、二甲苯、乙炔、苯酚、甲醛、乙二醇等等。

合成高分子是由单体聚合成的，所以高分子又叫高聚物。
合成高聚物的两种基本反应是加成聚合反应（简称加聚反应）和缩合聚合反应（简称缩聚反应）。
这在前面已经学过了。
在这里再结合实际事例作一些阐述。

\subsection{加聚反应}
不饱和烃(包括它们的衍生物)的单体通常是通过加聚反应聚合成高分子的。这里只介绍乙烯和丙烯的聚合。

\subsubsection{乙烯的聚合}
在不同温度、压强和有催化剂的条件下，乙烯分子作为单体能够发生加聚反应，生成性能不同的聚乙烯。

这个反应可以用下式来表示：
\[
\schemestart
\ce{$n$\,CH2=CH2}\arrow(.mid east--.mid west){->[\scriptsize 催化剂]}
\chemname{\chemfig{\vphantom{CH_2}-[@{op,.6}]CH_2CH_2-[@{cl,0.4}]}
\polymerdelim[delimiters ={[]},height=2pt,indice =\!n]{op}{cl}}{\footnotesize 聚乙烯}
\schemestop\chemnameinit{}
\]

当用氧气为催化剂、\qty{700}{atm} 和 \qty{150}{\celsius} 左右条件下，合成的聚乙烯在分子链上有较多支链，分子排列不整齐、密度较小（\qtyrange{0.91}{0.93}{g/cm^3}），强度也较低。
当用 \ce{TiCl4} 等制成的络合物作为催化剂、常压和低于 \qty{100}{ \celsius} 的条件下，合成的聚乙烯在分子链上只有很少支链，分子排列较整齐，密度较大（\qtyrange{0.94}{0.97}{g/cm^3}），强度也较大。
由此可见，改进催化剂，可以降低温度、压强，并合成性能较好的材料。

\subsubsection{丙烯的聚合}

二十世纪五十年代发现了用钛的氯化物（如 \ce{TiCl3}）等制成的络合物作为催化剂，能够使丙烯在聚合时生成的线型聚丙烯链作有规则的排列，即链上的甲基完全相同地排在一边 （\cref{fig:5-2}）。
\begin{figure}
  \begin{minipage}{\linewidth}\centering
    \chemfig[atom sep=3.5em]{-[:-20,0.5]C(-[:60,0.7]H)(-[:-60,0.7]H)-[:20]C(-[:110,0.7,,1]CH_3)(-[:-110,0.7]H)-[:-20]C(-[:60,0.7]H)(-[:-60,0.7]H)-[:20]C(-[:110,0.7,,1]CH_3)(-[:-110,0.7]H)-[:-20]C(-[:60,0.7]H)(-[:-60,0.7]H)-[:20]C(-[:110,0.7,,1]CH_3)(-[:-110,0.7]H)-[:-20]C(-[:60,0.7]H)(-[:-60,0.7]H)-[:20]C(-[:110,0.7,,1]CH_3)(-[:-110,0.7]H)-[:-20]C(-[:60,0.7]H)(-[:-60,0.7]H)-[:20]C(-[:110,0.7,,1]CH_3)(-[:-110,0.7]H)-[:-20,0.5]}
    \subcaption{聚丙烯链段示意图}\label{fig:5-2a}
  \end{minipage}
  \begin{minipage}{\linewidth}\centering
    \includegraphics{5-2b.pdf}
    \subcaption{聚丙烯链段球棍模型}\label{fig:5-2b}
  \end{minipage}
  \caption{聚丙烯结构示意图}\label{fig:5-2}
\end{figure}
丙烯的加聚反应可以简单表示如下：
\[
\schemestart
\ce{$n$\,CH3CH=CH2}\arrow(.mid east--.mid west){->[\scriptsize 催化剂]}
\chemfig{\vphantom{CH_2}-[@{op,.6}]CH(-[:-90]CH_3)-CH_2-[@{cl,0.4}]}
\polymerdelim[delimiters ={[]},height=2pt,indice =\!n]{op}{cl}
\schemestop
\]

这样聚合的聚丙烯链排列很整齐，容易达到较高程度结晶。
可以纺成很细而坚韧的丝，作纤维用；也可以制成薄膜，作为包装材料或很坚牢的带子。

\subsection{缩聚反应}
氨基酸聚合成蛋白质是缩聚反应的一个例子。
在这里，只叙述聚酯纤维和酚醛树脂的缩聚反应。

\subsubsection{合成聚酯纤维的缩聚反应}
聚酯纤维通常以对-苯二甲酸和乙二醇为单体，进行缩聚而成的。对-苯二甲酸和乙二醇的反应也是酯化反应。例如：
\begin{multline*} 
  \schemestart
  \chemname{\chemfig{HOOC-*6(-=-(-COOH)=-=) }}{\footnotesize 对-苯二甲酸}
  \+ \chemname{\ce{HOCH2CH2OH}}{\footnotesize 乙二醇} \arrow(.mid east--.mid west)
  \schemestop\\ 
  \chemfig{HOOC-*6(-=-(-C(=[:90]O)-OCH_2CH_2OH)=-=) } + \ce{H2O}
\end{multline*}

由于反应的生成物仍具有一个羧基和一个羟基两个能起反应的官能团，还能够跟对一苯二甲酸和乙二醇起反应，直到生成高分子化合物。
这个反应可以综合地表示如下：
\begin{multline*}
  \schemestart
  $n$\,\chemfig{HOOC-*6(-=-(-COOH)=-=)} \+ $n$\,\ce{HOCH2CH2OH}
  \arrow(.mid east--.mid west)
  \schemestop\\ 
  \schemestart
  \chemfig{\vphantom{CH_2}-[@{op,.6}]C(=[:90]O)-*6(-=-(-C(=[:90]O)-O-CH_2CH_2-O-[@{cl,0.4}])=-=)}
  \polymerdelim[delimiters ={[]},height=2pt,indice =\!n]{op}{cl}
\+ \ce{$2n$\,H2O}
  \schemestop
\end{multline*}

在上面的反应里，除生成高分子化合物外，同时生成小分子的水。
这是缩聚反应的特征，缩聚反应通常又叫逐步聚合反应。
在上述反应物所发生的缩聚反应里，生成的是线型高聚物。

\subsubsection{酚醛树脂}
酚醛树脂通常是由苯酚跟甲醛起反应而制得的。
苯酚在邻、对位都容易起反应，它能够跟甲醛形成线型的或体型的高聚物，在酸性条件下，苯酚过量时，苯酚主要在两个邻位起反应，合成线型酚醛树脂。
\[\schemestart
$n$\,\chemfig{[:-60]*6(=-=-=(-OH)-)} \+ $n$\,\chemfig{H-C(=[:90]O)-H} 
\arrow(.mid east--.mid west){->[\scriptsize 催化剂]}
\chemfig{\vphantom{CH_2}-[@{op,.6}]*6([:-30]-=-=(-[:0]CH_2-[@{cl,0.4}::0])-(-[:90]OH)=)}
  \polymerdelim[delimiters ={[]},height=2pt,indice =\!n]{op}{cl}
\+ $n$\,\ce{H2O}
\schemestop\]

这个反应还可以用下式来表示：
\[\schemestart 
\ce{ $n$C6H5OH } \+ \ce{ $n$HCHO } 
\arrow(.mid east--.mid west){->[\scriptsize 催化剂]}
\chemfig{\vphantom{CH_2}-[@{op,.6}]C_6H_3OHCH_2-[@{cl,0.4}]}
\polymerdelim[delimiters ={[]},height=2pt,indice =\!n]{op}{cl}
\+ \ce{ $n$H2O }
\schemestop\]

在碱性条件下，甲醛过量时，苯酚的苯环不仅在邻位而且在对位上都跟甲醛起反应，而得到体型的酚醛树脂（\cref{fig:5-3}）。
\begin{Experiment}
在试管里放 \qty{2.5}{g} 苯酚\footnote{上述实验可以用间苯二酚代替苯酚，缩聚反应可以进行得更迅速，效果更显著。操作手续相同。所用盐酸可略稀，即 10\% \ce{HCl}。}和 \qty{2.5}{ml} 40\% 甲醛溶液，再加入 \qty{1}{ml} 的浓盐酸。在沸水浴中加热。注意反应比较剧烈。发现溶液显浑浊，即可从水浴中取出试管。在另一个试管里加入 \qty{2.5}{g} 苯酚和 \qtyrange{3}{4}{ml} 的 40\% 甲醛溶液。在沸水浴中加热。再加入 \qty{1}{ml} 浓氨水。观察生成的树脂的颜色等。
\end{Experiment}
\begin{figure}
  \includegraphics{5-3.pdf}
  \caption{体型酚醛树脂结构示意图}\label{fig:5-3}
\end{figure}

酚醛树脂是一类合成树脂，它的单体，属于酚类的可以是苯酚、甲酚、间苯二酚等，属于醛类的可以是甲醛、糠醛等。
工业上最重要的酚醛树脂是由苯酚和甲醛缩聚而成的。

\begin{Practice}[习题]
\begin{question}
  \item 写出由单体苯乙烯（\ce{C6H5CH=CH2}）聚合成为聚苯乙烯的化学方程式。你认为这个反应是加聚反应还是缩聚反应？为什么？
  \item 根据你学过的知识，你认为加聚反应和缩聚反应有什么主要的区别？
  \item 缩聚反应是否一定要通过两种不同的单体才能发生？
\end{question}
\end{Practice}


\section{合成材料}
我们通常使用的合成高分子材料，即合成材料，主要有塑料、合成纤维和合成橡胶等。
经过缩聚反应而合成的高聚物主要作为塑料和合成纤维；经过加聚反应而合成的高聚物主要作为橡胶、塑料和合成纤维。
以下对它们作一些简要的介绍。

\subsection{塑料}
塑料是一类可塑性材料，主要成分是合成树脂。
学过的聚乙烯、聚丙烯、聚氯乙烯等都是生产塑料的合成树脂。
好些合成树脂是属于线型结构的，具有热塑性，用它制成的塑料就是\Concept{热塑性塑料}。
相反地，象体型的酚醛树脂，具有热固性，用它制成的塑料就是\Concept{热固性塑料}。

生产塑料制品的原料，除合成树脂外，一般还要加入能改进机械性能的增强材料和填料，能提高塑性的增塑剂，防老化的防老剂\footnote{增强材料有碳纤维、硼纤维、玻璃纤维、棉、麻等；填料有木粉、石棉、云母、粘土等；增塑剂，如邻苯二甲酸二辛酯等，用于增进聚氯乙烯的可塑性和柔软性；防老剂种类很多，有的使塑料在热加工时稳定，有的能抗氧气的侵蚀，等等。}等等添加剂，经压制成型或挤压成型等加工处理，才能制成塑料成品。

生产塑料的合成树脂种类较多，除前面已提到的几种外，还有聚四氟乙烯、有机玻璃、环氧树脂等。

以下简单介绍聚氯乙烯的制法、性质、用途等。
聚氯乙烯是由单体氯乙烯（\ce{CH2=CHCl}）经加聚反应而合成的。
反应可以简单表示如下：
\[\schemestart 
\ce{ $n$CH2=CHCl }  
\arrow(.mid east--.mid west)
\chemfig{\vphantom{CH_2}-[@{op,.6}]CH_2CHCl-[@{cl,0.4}]}
\polymerdelim[delimiters ={[]},height=2pt,indice =\!n]{op}{cl}
\schemestop\]

聚氯乙烯是线型高聚物，在长链的碳原子上间隔地出现极性基氯原子，使主链较难转动，同时也增加了高分子链间的引力。
所以聚氯乙烯较为强韧，能耐稀酸，具有热塑性，但不耐热。
要加入增塑剂才能使聚氯乙烯制品成为质软而有弹性的软质聚氯乙烯，它大量地被用制农用薄膜、电线和电缆的包皮、软管和日常生活用品等。
在聚氯乙烯树脂中不加或少加增塑剂就制成硬质聚氯乙烯，它的特点是硬而轻、强度好，可作建筑材料、绝缘材料、耐腐蚀材料等。

聚氯乙烯受热易发生消去反应而放出氯化氢。
聚氯乙烯这种遇热易分解的性质，使得在加工过程中增加困难，为了不致放出氯化氢，必须加入能吸收它的碱性物质作为稳定剂。

聚氯乙烯塑料在使用过程中，如果不注意环境的影响，容易发生老化。
如受日光的曝晒，受空气、水以及腐蚀性气体的作用，受辐射，受热等因素以及增塑剂的挥发，时间久了变硬发脆，开裂，长霉等。

塑料的品种比较多，贮存、使用条件也常不同，因而有各种不同的老化特征。
如农用薄膜经日晒雨淋，发生变色、变脆、透明度下降等；户外电缆、电线包皮使用日久变硬破裂等；有机玻璃的手表表面用久后透明度下降等。

此外，随着生产的现代化和科学技术的迅速发展，根据需要制出了许多特殊用途的塑料，如工程塑料、增强塑料、功能高分子等等。

\begin{Reading}[补充阅读]{}
\subparagraph*{工程塑料} 聚酰胺 1010
\[ 
\chemfig{\vphantom{CH_2}-[@{op,.6}]N-[,1.5]{(CH_2)_{10}}-[,1.6]N(-[:90]H)-C(=[:90]O)-[,1.5]{(CH_2)_{8}}-[,1.6]C(=[:90]O)-[@{cl,0.4}]}
\polymerdelim[delimiters ={[]},indice =\!n]{op}{cl}
\]
是一种工程塑料，它是由癸二胺（\ce{H2N(CH2)10NH2}）跟癸二酸（\ce{HOOC(CH2)8COOH}）缩聚而成的。
用它制造的机械零件，如齿轮、轴承等具有摩擦力小、自动润滑、使用寿命长、节省动力等特点。
聚酰胺~6 等也可以作为工程塑料。
由丙烯腈、 丁二烯、苯乙烯共同聚合而成的 ABS 工程塑料是发展很迅速的一种。

\subparagraph*{增强塑料} 
用碳纤维、硼纤维、玻璃纤维增强的酚醛塑料可用于制造宇宙飞船、人造卫星、洲际导弹的部件。
环氧树脂强度较差，当它跟玻璃纤维复合，就制得强度类似钢的增强塑料——玻璃钢。
用碳纤维增强的环氧树脂，强度超过最强的钢。

\subparagraph*{功能高分子} 
通过化学反应把适当的官能团引入高聚物中，使生成的高聚物具有离子交换、光敏等等功能的是功能高分子。
例如，在硬水软化中应用的离子交换树脂即是在高聚物中引入酸或碱的官能团，能起交换离子的功能；光敏性高分子有在光照射时发生快速反应的官能团，光敏性高分子用于印刷线路等的制作。
\end{Reading}

\subsection{合成纤维}
棉花、羊毛、木材和草类的纤维都是天然纤维。
粘液丝纤维是天然纤维经人们加工而成的人造纤维。
利用石油、天然气、煤和农副产品作原料制成单体，经聚合反应制成的是合成纤维。
合成纤维和人造纤维又统称化学纤维。

合成纤维是在二十世纪三十年代开始生产的，它为人民生活提供了各种经久耐用而美观的衣着材料。
合成纤维是由线型高分子加工制成的。
高分子链之间以分子间力相作用，有的还有氢键的作用，有很好的强度、弹性，耐磨、耐腐蚀性也较好。
合成纤维品种繁多，如聚酯纤维、聚酰胺纤维、聚丙烯腈纤维、聚烯烃纤维等。
以下对聚酰胺等纤维作简单的介绍。

\subparagraph*{聚酰胺纤维}
聚酰胺高分子的链上有酰胺基
\[ \chemfig{-C(=[:90]O)-NH-}\]

聚酰胺纤维有多种，如前面讲过的工程塑料聚酰胺~1010 就是其中之一，其它还有聚酰胺~6 （即锦纶）、聚酰胺~66 等。
这里只介绍聚酰胺~6 的生产。
生产聚酰胺~6 的原料有好多种，可以从苯、环己烷或苯酚为起始的原料，经过几步的反应，制成单体己内酰胺（\rule[-2.5ex]{0pt}{6ex}\schemestart H\OX{n,N}(CH$_2$)$_5$\OX{c,C}O\redox(n,c)[][-1]{}\schemestop），它是含有 6 个碳原子的酰胺，并连接成环状，所以叫己内酰胺。
纯的己内酰胺不会自动发生聚合，当在高温和引发剂（通常是水）的作用下，己内酰胺被水解，环被打开。
\[ \schemestart
H\OX{n,N}(CH$_2$)$_5$\OX{c,C}O\redox(n,c)[][-1]{} \+ \ce{H2O}
\arrow(.mid east--.mid west)
\chemfig{H_2N{(CH_2)_5}-[,1.2,,,draw=none]C(=[:90]O)-OH}
\schemestop
\]

生成的氨基己酸有两个官能团，一个氨基己酸的分子的氨基能够跟另一个分子的羧基发生反应。
\begin{multline*}
  \schemestart
  \chemfig{H_2N{(CH_2)_5}-[,1.2,,,draw=none]C(=[:90]O)-OH}
  \+
  \chemfig{H_2N{(CH_2)_5}-[,1.2,,,draw=none]C(=[:90]O)-OH}
  \arrow(.mid east--.mid west)\schemestop\\
  \schemestart
  \chemfig{H_2N{(CH_2)_5}-[,1.2,,,draw=none]C(=[:90]O)-NH{(CH_2)_5}-[,1.2,,,draw=none]C(=[:90]O)-OH}
  \+ \ce{H2O}
\schemestop
\end{multline*}

生成物进一步发生缩聚反应，生成聚酰胺~6，总的化学方程式可以表示如下：
\[\schemestart
$n$\,\chemfig{H_2N{(CH_2)_5}-[,1.2,,,draw=none]COH(=[:90]O)}
\arrow(.mid east--.mid west)
\chemfig{\vphantom{CH_2}-[@{op,.6}]N(-[:90]H)-[,,,,draw=none]{(CH_2)_{5}}-[,,,,draw=none]C(=[:90]O)-[@{cl,0.4}]}
\polymerdelim[delimiters ={[]},height=2pt,indice =\!n]{op}{cl}
\+ \ce{$n$H2O}
\schemestop\]

聚酰胺纤维是目前合成纤维中强度最大的，用于制降落伞、轮胎帘子线、渔网等。
为什么聚酰胺纤维有这样优良性能呢？
因为它的分子链之间除分子间力外，还有氢键的作用（\cref{fig:5-4}），加强了分子间的引力，有利于形成排列比较整齐的结晶区域（\cref{fig:5-5}），再经拉伸，结晶区域排列更整齐，分子间力和氢键的作用发挥得更好。
由于结晶区域跟非结晶区域结合在一起，使聚酰胺纤维具有优良的性能。
\begin{figure}
  \begin{minipage}[b]{0.48\linewidth}\centering
    \chemfig[atom sep=2.3em]{-[,0.7,,,dotted,very thick]H-N(-[:-60]CH_2-[:-120]\phantom{C})-[:60]C(-[:120]CH_2-[:60]\phantom{C})=O-[,0.9,,,dotted,very thick]H-N(-[:-60]CH_2-[:-120]\phantom{C})-[:60]C(-[:120]CH_2-[:60]\phantom{C})=O-[,0.7,,,dotted,very thick]}
    \caption{聚酰胺链间的氢键示意图}\label{fig:5-4}
  \end{minipage}
  \begin{minipage}[b]{0.48\linewidth}\centering
    \includegraphics{5-5.pdf}
    \caption{高分子材料部分结晶示意图}\label{fig:5-5}
  \end{minipage}
\end{figure}

\begin{Reading}[补充阅读]{}
\subparagraph*{聚酯纤维} 聚酯纤维的链里含有酯基
\[\chemfig{-C(=[:90]O)-O-}\]
酯基也是极性基团，使高分子链间有着强的分子间力，增加了在拉伸时获得高度结晶的可能性。
所以聚酯纤维是一种性能比较优良的合成纤维，它的抗折皱性特别好，弹性、耐磨性也比较好。对苯二甲酸（\chemfig[atom sep=0.8em]{HOOC-*6(-=-(-COOH)=-=)}）可以跟乙二醇（\ce{CH2OHCH2OH}）起缩聚反应而合成俗名涤纶的聚酯纤维，
\[
\chemfig{\vphantom{CH_2}-[@{op,.6}]OCH_2CH_2OCO-*6(-=-(-CO-[@{cl,0.4}])=-=)}
\polymerdelim[delimiters ={[]},height=2pt,indice =\!n]{op}{cl}
\]
即的确良。


\subparagraph*{聚丙烯腈纤维}
\[
\chemfig{\vphantom{CH_2}-[@{op,.6}]CH(-[:-90]CN)-CH_2-[@{cl,0.4}]}
\polymerdelim[delimiters ={[]},hight=3pt,depth=5pt,indice = \!n]{op}{cl}
\]
因分子链上有氰基 \ce{-CN}，虽然分子链不整齐，但有氰基极性基团，具有较大分子间力，也可以制成纤维。
聚丙烯腈纤维强度不足，但弹性好，通常叫人造羊毛。


\subparagraph*{聚烯烃纤维} 
聚烯烃纤维主要是聚乙烯纤维和聚丙烯纤维。
\[\schemestart
\chemname{\chemfig{\vphantom{CH_2}-[@{op,.6}]CH_2CH_2-[@{cl,0.4}]}
\polymerdelim[delimiters ={[]},indice = \!n]{op}{cl}}{\footnotesize 聚乙烯纤维}\qquad
\chemname{\chemfig{\vphantom{CH_2}-[@{op,.6}]CHCH_3CH_2-[@{cl,0.4}]}
\polymerdelim[delimiters ={[]},indice = \!n]{op}{cl}}{\footnotesize 聚丙烯纤维}
\schemestop\chemnameinit{}\]
烯烃当用络合物作催化剂起加聚反应时，生成具有整齐排列的高聚物可以作为合成纤维使用。
例如聚乙烯链上只有氢原子，没有别的基团，分子排列容易整齐，拉伸时结晶程度较高，能作纤维用。
\end{Reading}

合成纤维在长久使用后发生老化，主要表现为强度下降、褪色、断裂等。
老化的原因主要是由于合成纤维结构上具有羰基等活泼的基团受阳光、氧气、臭氧、热、水、高能辐射、工业气体（如 \ce{CO2}、\ce{NH3}、\ce{HCl}）、霉菌等外界因素的作用。
使用时如果注意避免这些外界因素会延缓合成纤维制品的老化。

\subsection{橡胶}
橡胶有天然橡胶和合成橡胶。
十八世纪末人们开始利用天然橡胶制作橡皮等文具用品。
十九世纪中叶天然橡胶才真正进入实用阶段。
合成橡胶在二十世纪初开始合成，四十年代起得到了迅速的发展。
天然橡胶在性能方面比较全面，如弹性、电绝缘性、加工性能等都比较好。
合成橡胶在某些性能上比较突出，如有的耐高温、耐低温，有的耐油，有的有很好的气密性，等等。
橡胶是制造飞机、军舰、拖拉机、收割机、汽车、水利排灌机械、医疗器械等所必需的材料。
日常生活中许多用品的生产也离不开橡胶。

\subsubsection{天然橡胶}
天然橡胶主要来源于橡胶树和橡胶草的胶乳。
胶乳的主要成分是聚异戊二烯。
由橡胶树（三叶橡胶树）和橡胶草的胶乳炼成的橡胶弹性好。
天然橡胶是聚异戊二烯，它的结构式如下：
\[ 
\chemfig{\vphantom{CH_2}-[@{op,.6}]CH_2-[:-60,,1,1]C(-[:-120,,1,1]CH_3)=C(-[:-60]H)-[:60,1.05]CH_2-[@{cl,0.4}]}
\polymerdelim[delimiters ={[]},hight=3pt,depth=50pt,indice = \!n]{op}{cl}
\]

制天然橡胶的胶乳是一种乳白色胶体，从橡胶树树皮内的乳管流出。
胶乳浓缩后可供生产橡胶制品，如医用手套等。
胶乳经加醋酸凝结，压片，干燥，制成生胶片；或经凝结、造粒、干燥，制成颗粒橡胶。
这种尚未炼制的生胶受热变软，遇冷变硬、发脆，不易成型，容易磨损；易溶于汽油、四氯化碳等有机溶剂；分子内具有双键易起加成反应，容易老化。
生产上要进行硫化等一系列加工过程，以改善橡胶制品的性能。
硫化是使线型高分子经交联而变成网状结构的体型高分子（\cref{fig:5-6}）。

\begin{figure}
  \begin{minipage}[b]{0.48\linewidth}\centering
    \includegraphics{5-6a.pdf}
    \subcaption{交联}\label{fig:5-6a}
  \end{minipage}
  \begin{minipage}[b]{0.48\linewidth}\centering
    \chemfig{-CH_2-C(-[:90]S-[:90]S-[:90]C(-[:90]CH_3)(-[:180]CH_2-[:180]\phantom{C})-CH(-[:90]\phantom{C})-CH_2-)(-[:-90]CH_3)-CH(-[:-90]\phantom{C})-CH_2-}
    \subcaption{交联的结构}\label{fig:5-6b}
  \end{minipage}
  \caption{橡胶硫化示意图}\label{fig:5-6}
\end{figure}

\begin{Reading}[补充阅读]{}
为什么需要交联呢？
生橡胶的线型分子链是柔软的，链段能够移动，在外力作用下由蜷曲变为伸展，外力消失后又能够恢复，这是具有弹性的必要条件。
但链跟链之间容易滑动，强度和韧性很差，必须经过交联，使分子链跟分子链之间有化学键的联系，才能增加强度和韧性。
交联的办法通常是硫化。
硫化过程是很复杂的，常是形成双硫键起交联作用。
硫化时所加硫通常占总量的 3\% 左右，交联程度比较低，得到的橡胶制品有很好的弹性、韧性等。
如果加硫达 30\% 左右，链上的 $\pi$ 键几乎都被交联，生成高度交联的体型高分子，就是硬橡皮。
\end{Reading}

\subsubsection{合成橡胶}
人们在长期的生产和科学实验中认识了天然橡胶的结构，得到启发，终于成功地合成了多种合成橡胶。
合成橡胶以石油、天然气为原料，生产的二烯烃和烯烃作为单体，如丁二烯、异戊二烯、氯丁二烯、苯乙烯、异丁烯、丙烯腈等，其中最重要的是丁二烯、通常应用的合成橡胶有丁苯橡胶、顺丁橡胶、氯丁橡胶等，它们都是通用橡胶。
\subparagraph*{丁苯橡胶}
丁苯橡胶是由两种单体，即丁二烯和苯乙烯在引发剂\footnote{引发剂是指能够引发反应的物质（反应物以外的），在这里是指引发聚合反应。}或催化剂作用下，起加聚反应而生成的。
反应过程很复杂，下面的化学方程式只是简单的表示方式。
\begin{multline*}\schemestart  
\ce{$n$CH2=CH-CH=CH2} 
\+ $n$
\chemfig{CH_2=CH-[:-90]*6(=-=-=-)}
\arrow(.mid east--.mid west)
\schemestop\\[5pt]
\chemfig{\vphantom{CH_2}-[@{op,.6}]CH_2-CH=CH-CH_2-CH_2-CH(-[:-90]*6(=-=-=-))-[@{cl,0.4}]}
\polymerdelim[delimiters ={[]},hight=3pt,depth=5pt,indice = \!n]{op}{cl}
\end{multline*}

丁苯橡胶的耐日光、耐气候和抗老化性比较好，它的链节里还有一个化学性质比较稳定的苯环，提高了它的热稳定性。
但丁苯橡胶的硫化速度慢、耐寒性差。

\subparagraph*{顺丁橡胶}
顺丁橡胶是当前发展十分迅速的一个合成橡胶品种。
一般的聚丁二烯性能并不好，但是用钛的络合物作为催化剂，以 1,3-丁二烯为单体聚合而成的顺丁橡胶，由于象天然橡胶那样 \ce{CH2} 基团都顺列在一边而结构比较有规则，因而也象天然橡胶那样具有良好的性能。
\[\schemestart 
\ce{$n$CH2=CH-CH=CH2} \arrow(.mid east--.mid west){->[\scriptsize 络合物催化剂]}
\chemname{
\chemfig{\vphantom{CH_2}-[@{op,.6}]CH_2-[:-60,1.2,1,1]CH=CH-[:60,1.2,1,1]CH_2-[@{cl,0.4}]}
\polymerdelim[delimiters ={[]},height=3pt,depth =25pt,indice = \!n]{op}{cl}
}{\footnotesize 顺丁橡胶}
\schemestop\chemnameinit{}\]

我国有丰富的稀土金属资源，用我国研制的稀土金属络合物作为催化剂，能获得纯度较高的顺丁橡胶。
顺丁橡胶的耐磨性高、弹性良好，广泛用于制造轮胎。

\begin{Reading}[补充阅读]{}
\subparagraph*{氯丁橡胶}是由氯丁二烯作为单体聚合而成的，结构简式是：
\[
\chemfig{\vphantom{CH_2}-[@{op,.6}]CH_2-C(-[:-90]Cl)=CH-CH_2-[@{cl,0.4}]}
\polymerdelim[delimiters ={[]},indice = \!n]{op}{cl}
\]
\end{Reading}


\subparagraph*{特种橡胶}有耐油很好的聚硫橡胶和耐严寒和高温的硅橡胶。
\[
\schemestart
\chemname{\chemfig{\vphantom{CH_2}-[@{op,.6}]CH_2CH_2S_4-[@{cl,0.4}]}
\polymerdelim[delimiters ={[]},height = 2pt, indice = \!n]{op}{cl}}{\footnotesize 聚硫橡胶}\qquad
\chemname{\chemfig{\vphantom{CH_2}-[@{op,.6}]Si(-[:90]CH_3)(-[:-90]CH_3)-O-[@{cl,0.4}]}
\polymerdelim[delimiters ={[]},height = 2pt, indice = \!n]{op}{cl}}{\footnotesize 硅橡胶}
\schemestop\chemnameinit{}
\]

一般橡胶制品只达到轻度交联，高分子链上还有双键，受氧气、臭氧，受日光，特别是高能辐射的作用，还容易老化。
贮存或使用日久的汽车轮胎、自行车胎会发生龟裂，乳胶手套会发粘或变硬脆裂，其它橡胶制品会出现弹性下降，变硬开裂或变软发粘等。

\begin{Theorem}{讨论}
  根据已经学过的知识，讨论聚氯乙烯、顺丁橡胶、聚酰胺~6 等合成材料的单体类别、聚合的反应类型、结构（线型或体型）各有什么共同的和不同的地方。
\end{Theorem}


\begin{Practice}
  \begin{question}
    \item 什么是热塑性塑料，什么是热固性塑料？这两类塑料在结构上有什么不同？
    \item 为什么有的塑料在制造过程中要加增塑剂？增塑剂的作用是什么？
    \item 人造纤维和合成纤维有什么区别？
    \item 为什么天然橡胶和合成橡胶都要经过硫化？
    \item 试比较塑料、合成纤维、橡胶等三类高分子材料的结构。
    \item 从塑料、合成纤维和橡胶中，各举两个例子来说明它们的单体、链节和反应的化学方程式。
  \end{question}
\end{Practice}

\section*{内容提要}
\addcontentsline{toc}{section}{内容提要}\setcounter{subsection}{0}
\subsection*{一、高分子的式量、结构和性质}
高分子的式量很大，高分子的式量是平均式量。

高分子的结构大体可以分为线型结构和体型结构两大类。
高分子链是由链节以一定聚合度联接而成的。
高分子材料的性质跟高分子链中化学键的键能大小和高分子链间的分子间力大小有直接关系。
高分子的式量大，链也长，分子间力也大，链间形成氢键的，作用力更大。

线型高分子能够溶解在适当的溶剂里，体型高分子不能溶解。
线型高分子受热会熔化，体型高分子受热不熔化。
高分子材料还有高强度、塑性、电绝缘性、耐腐蚀性、密度小等优良性能。
不同的高分子材料各有它们突出的性质。

\subsection*{二、合成材料}
\begin{enumerate}
  \item 塑料\quad 由线型高分子制成的是热塑性塑料，体型高分子制成的是热固性塑料。
  \item 合成纤维\quad 合成纤维是线型高分子加工制成的，有很好的强度。
  \item 橡胶\quad 线型的橡胶高分子链，通常要经过硫化，成为网状的高分子材料。硫化橡胶一般交联程度少，富于弹性。根据使用不同，橡胶分为通用橡胶和特种橡胶。
\end{enumerate}

\begin{Review}
\begin{question}
  \item 下列说法是否正确？说明理由。
  \begin{tasks}[label=(\arabic*)]
    \task 合成橡胶的单体都是二烯烃。
    \task 高分子化合物的溶液既是单个分子分散在溶剂里的真溶液，又是具有丁达尔现象等的胶体溶液。
    \task 合成纤维都是由单体经缩聚反应而合成的。
  \end{tasks}
  \item 高分子的式量跟小分子的式量有什么不同？为什么说高分子要具有相当大的式量才具有高分子物质应有的特性？
  \item 推测下列高聚物的单体和写出由单体聚合成高聚物的简单的化学方程式。
  \begin{tasks}[before-skip=10pt,after-skip=10pt,after-item-skip=7pt](2)
    \task \chemfig{\vphantom{CH_2}-[@{op,.6}]CH_2-CH(-[:-90]F)-[@{cl,0.4}]}
\polymerdelim[delimiters ={[]},height = 2pt, indice = \!\!n]{op}{cl}
    \task \chemfig{\vphantom{CH_2}-[@{op,.6}]CH_2-CH(-[:-90]CH_3)(-[:90]CH_3)-[@{cl,0.4}]}
\polymerdelim[delimiters ={[]},height = 2pt, indice = \!\!n]{op}{cl}
    \task \chemfig{\vphantom{CH_2}-[@{op,.6}]C(-[:-90]F)(-[:90]F)-C(-[:-90]F)(-[:90]F)-[@{cl,0.4}]}
\polymerdelim[delimiters ={[]},height = 2pt, indice = \!\!n]{op}{cl}
  \end{tasks}
  \item 试说明人们从发现三叶橡胶树的胶乳到合成高分子链排列整齐的顺丁橡胶，经过了哪些具有重要意义的突破过程。
  \item 根据你所了解的知识，说明合成材料对工农业生产、国防建设和日常生活的重要意义。
\end{question}
\end{Review}

\begin{Exercise}*[总复习题]
\begin{question}
  \item 以铁为例，说明过渡元素在原子结构和性质上不同于主族元素的一些特点。
  \item 怎样用化学方法来除去
  \begin{tasks}(2)
    \task 铜粉中混有的铁粉，
    \task \ce{FeCl2} 中混有的 \ce{CuCl2}，
    \task \ce{FeCl2} 中混有的 \ce{FeCl3}，
    \task 铁粉中混有的铝粉。
  \end{tasks}
  \item 现有一小块铜银合金，试设计分离、鉴别的实验步骤，写出反应的化学方程式。
  \item 在实验室怎样从 \ce{FeSO4} 制取 \ce{FeCl3}，写出反应的各步化学方程式。
  \item A 是一种黑色固体，它跟盐酸反应生成无色气体 B，B 能使石蕊试纸变红，反应后溶液变成淡绿色，在这溶液中加入 \ce{NaOH} 溶液，生成白色沉淀 C，C 在空气中迅速变成灰绿色，最后变成红褐色沉淀 D，试判断 A、B、C、D 各是什么物质。
  \item 写出下列反应的离子方程式：
  \begin{tasks}
    \task \ce{AgNO3} 溶液跟氨水，
    \task \ce{FeCl3} 溶液跟 \ce{KSCN} 溶液，
    \task \ce{FeCl3} 溶液跟 \ce{C6H5OH} 溶液，
    \task \ce{CuSO4} 溶液跟氨水，
    \task \ce{Au} 跟 \ce{KCN} 溶液。
  \end{tasks}
  \item 现有 $0.1M$ \ce{CoCl2.5NH3} 溶液 \qty{200}{ml}，加足量 \ce{AgNO3} 溶液后，生成 \ce{AgCl} 沉淀 \qty{5.74}{g}。
  
  这种络合物的络离子是\CJKunderline[hidden]{\hspace{1em}\ce{[Co(NH3)5Cl]}\hspace{1em}}，中心离子是\CJKunderline[hidden]{\hspace{1em}\ce{Co^2+}\hspace{1em}}。配位体是\CJKunderline*[hidden]{\hspace{1em}\ce{NH3}、\ce{Cl-}\hspace{1em}}，配位数是\CJKunderline[hidden]{\hspace{1em}6\hspace{1em}}。
  \item 根据下述条件，推断几种烃的化学式。
  \begin{tasks}
    \task 某烃 \qty{0.1}{mol} 完全燃烧后，生成 \ce{CO2} \qty{0.3}{mol} 和 \ce{H2O} \qty{0.3}{mol}。
    \task 在标准状况下，\qty{20}{ml} 某气态烃完全燃烧后，生成 \qty{40}{ml} \ce{CO2} 和 \qty{60.1}{ml} \ce{H2O}（气态）。
    \task 某气态烯烃完全燃烧时，消耗氧气的体积是它本身体积的 6 倍。
  \end{tasks}
  \item 写出具有下列碳链的炔烃的结构式，并写出它们的名称：
  \begin{tasks}[before-skip=10pt,after-skip=10pt,after-item-skip=7pt](2)
    \task \chemfig{C-C(-[:-90]C)-C-C-C(-[:-90]C)-C}
    \task \chemfig{C-C(-[:-90]C)-C-C(-[:-90]C)-C-C}
    \task \chemfig{C-C(-[:-90]C)-C(-[:-90]C-[:-90]C)-C}
  \end{tasks}
  \item 根据系统命名法注明下列各种物质的名称：
  \begin{tasks}[before-skip=10pt,after-skip=10pt,after-item-skip=7pt](2)
    \task \chemfig{CH_3-C(-[:90]Cl)(-[:-90]CH_3)-CH=CH_2}
    \task \chemfig{CH_3-CH(-[:90]CH_3)-CH(-[:-90]CH_2-[:-90]OH)-CH_3}
    \task \chemfig{CH_3-CH=CH-C(-[:90]CH_3)=CH_2}
    \task \chemfig{*6(=(-[:-90]C_2H_5)-=(-[:0]C_2H_5)-=-)}
    \task \chemfig{CH_3-CH(-[:90]CH_3)-C(=[:90]O)-H}
    \task \chemfig{\vphantom{CH_2}-[@{op,.6}]CH(-[:-90]CH_3)-CH_2-[@{cl,0.4}]}\polymerdelim[delimiters ={[]},indice=\!n]{op}{cl}
  \end{tasks}
  \item 写出三种制取溴乙烷的方法。
  \item 下列各种物质中，互为同分异构体的是\CJKunderline[hidden]{\hspace{1em}E 和 H\hspace{1em}}，互为同系物的是\CJKunderline*[hidden]{\hspace{1em}D 和 G，C 和 F\hspace{1em}}。
  \begin{tasks}[label=\Alph*.](3)
    \task \ce{C2H4}
    \task \ce{C2H2}
    \task \ce{HCHO}
    \task \ce{C6H6}
    \task \ce{CH3COOC2H5}
    \task \ce{C2H5CHO}
    \task \ce{C6H5C2H5}
    \task \ce{C2H5COOCH}
  \end{tasks}
  \item 怎样用化学方法来除去
  \begin{tasks}(2)
    \task 苯中混有的苯胺，
    \task 苯中混有的苯酚，
    \task 乙烷中混有的乙烯，
    \task 乙酸乙酯中混有的乙酸。
  \end{tasks}
  \item 填空：
  \begin{tasks}
    \task 某学生用乙醇跟浓硫酸加热来制取乙醚，结果得到的产物主要是乙烯，这主要是由于\CJKunderline*[hidden]{\hspace*{1em}温度过高，乙醇分子脱水生成乙烯\hspace*{1em}}。
    \task 某学生试用 1-溴丙烷跟氢氧化钠的乙醇溶液起反应来制取丙醇，结果得到的产物主要是丙烯，这主要是由于\CJKunderline*[hidden]{\hspace*{1em}1-溴丙烷跟\ce{NaOH} 的水溶液发生反应，生成丙醇；如与 \ce{NaOH} 的乙醇溶液共热，则发生消去反应生成丙烯 \hspace*{1em}}。
  \end{tasks}
  \item 某有机物的化学式是 \ce{C3H4O2}，它的水溶液显酸性，能跟 \ce{Na2CO3} 溶液起反应，又能使溴水褪色。写出这种有机物的结构式。
  \item 某种烃含 \ce{C} 92.3\%，含 \ce{H} 7.7\%。在标准状况下，\qty{1}{L} 这种气体的质量是 \qty{3.49}{g}，求这种烃的化学式。
  \item 某气态化合物中各元素的质量比为 $\ce{C}:\ce{H}:\ce{O}=6:1:8$。在同温同压下，这种气体的质量是同体积氧气质量的 0.94 倍，求这种化合物的化学式。
  \item 今有 $A$、$B$ 两种含碳、氢、氧三种元素的链烃衍生物，它们具有相同的百分组成，$A$ 的式量是 44，$B$ 的式量是 88。
  
  $A$ 能发生银镜反应而 $B$ 不能。把 $A$ 加入新制的 \ce{Cu(OH)2} 并加热到沸腾，有红色沉淀产生，并生成有机物 $C$。$B$ 在无机酸催化作用下水解生成 $C$ 和 $D$，$D$ 在有催化剂（\ce{Cu} 或 \ce{Ag}）存在条件下加热，能够被空气氧化生成 $A$，$C$ 的钠盐（无水）和碱石灰混和后加热可得到一种式量为 16 的气态烷烃。

  写出 $A$、$B$、$C$、$D$ 的名称和结构简式，写出上述各步反应的化学方程式。
  \item 脂肪、淀粉、蛋白质是三种重要的营养成分，其中\CJKunderline[hidden]{\hspace{1em}脂肪\hspace{1em}}不是高分子化合物，脂肪的水解产物是\CJKunderline[hidden]{\hspace{1em}高级脂肪酸\hspace{1em}}和\CJKunderline[hidden]{\hspace{1em}甘油\hspace{1em}}，淀粉的水解产物是\CJKunderline[hidden]{\hspace{1em}葡萄糖\hspace{1em}}，蛋白质的水解产物是\CJKunderline[hidden]{\hspace{1em}$\alpha$-氨基酸\hspace{1em}}。
  \item 将正确答案的标号（A、B、C、D）填在括弧内。
  \begin{enumerate}[(1),itemindent=2.4em]
    \item 苯酚和甲醛起反应生成酚醛树脂，发生的是\hfill（\qquad ）
    \begin{tasks}[label=\Alph*](4)
      \task 加成反应，
      \task 缩聚反应，
      \task 酯化反应，
      \task 加聚反应。
    \end{tasks}
    \item 淀粉和纤维素是\hfill（\qquad ）
    \begin{tasks}[label=\Alph*]
      \task 同分异构体，
      \task 同系物，
      \task 都是葡萄糖的加聚产物，
      \task 组成相同、式量不同的天然高分子化合物。
    \end{tasks}
  \end{enumerate}
\end{question}
\end{Exercise}
